\documentclass{article}

\usepackage{amsmath}
\usepackage[style=apa,natbib=true]{biblatex}
\addbibresource{thesis.bib}

\begin{document}

\section{Binomické stromy}\label{sec:binom}

Nemůžeme s jistotou odhadnout, jak se bude vyvíjet cena akcie. Jedním ze způsobů, jak se s tímto vypořádat, je modelovat očekávanou hodnotu zvětšení či zmenšení hodnoty akcie. Na tomto principu funguje oceňování opce pomocí binomických stromů.

Binomický strom je strom, jehož každý vrchol má nejvíce právě jednoho rodiče a maximálně dva potomky. %TODO: Zdroj
Tento strom využijeme k zjednodušení vývoje ceny akcie v čase: Místo toho, aby akcie mohla nabýt libovolné hodnoty z reálných čísel, řekneme, že cena akcie se multiplikativně zvedne o specifickou hodnotu, pokud její cena roste, a cena akcie se mulitplikativně sníží o specifickou hodnotu, pokud cena klesne. Je-li počáteční cena akcie $S$, můžeme tyto násobující konstanty označit jako $u$ pro pohyb nahoru, a $d$ pro pohyb dolů, a cenu po růstu $uS$ a cenu po snížení $dS$. Všechny tyto zmíněné ceny akcií jsou popsány v binomickém stromu, kde $S$ je rodič a $uS, dS$ jsou potomci. Obvykle volíme $u > 1$ a $d < 1$, aby $u$ reprezentovalo růst a $d$ pokles ceny; přesné hodnoty závisí na zvoleném modelu. %TODO: Zdroj

%TODO: Obrázek stromu

Rádi bychom dokázali, že ceny opcí na akcích jsou jednoznačně určitelné. K tomuto účelu, v rámci tohoto zjednodušeného modelu vývoje ceny akcie můžeme vytvořit dvě portfolia: 

\begin{description}
  \item[Opční portfolio] skládající se z jedné call opce o ceně $C$.
  \item[Syntetické portfolio] skládající se z $\Delta$ podílů akcie (stejné jako call opce) a z hodnoty investice do bezrizikového aktiva $\beta$.
\end{description}

$\Delta$ a $\beta$ nastavíme tak, aby při obou možnostech vývoje ceny akcie se hodnota obou portfolií rovnala. Čili:

\begin{gather}
  \Delta u \, S e^{qT} + \beta \, e^{rT} = C_u, \\ 
  \Delta d \, S e^{qT} + \beta \, e^{rT} = C_d,
\end{gather}

kde $C_u$ vyjadřuje cenu opce při růstu ceny akcie, $C_d$ vyjadřuje cenu opce při poklesu ceny akcie a $T$ vyjadřuje délku periody. Výdělek z akcie zvyšujeme o $e^{qT}$, což je výše vyplácené dividendy, o které očekáváme, že ji reinvestujeme, a půjčené peníze musíme diskontovat o $e^{rT}$. \citep{CoxRossRubinstein1979} \citep[Kapitola~13]{McDonald2013}

Odečtením rovnic a následnou úpravou získáme: 

\begin{equation}
  \Delta = e^{-qT} \left( \frac{C_u - C_d}{S(u - d)} \right),
\end{equation}

Následně osamostatněním $\Delta$, opět odečtením rovnic a následnou úpravou vyplyne:

\begin{equation}
  \beta = e^{-rT} \left( \frac{u C_d - d C_u}{u - d} \right).
\end{equation}

Jelikož jsme nastavili obě portfolia tak, aby se rovnala, jsme schopni vyjádřit cenu opce jako:

\begin{equation}
  \label{eq_value}
  \Delta S + \beta = e^{-rT} \left( C_u \frac{e^{(r - q) T} -d}{u - d} + C_d \frac{u - e^{(r - q) T}}{u - d} \right)
\end{equation}

\citep{CoxRossRubinstein1979} \citep[Kapitola~13]{McDonald2013}

Očekáváme, že $u > e^{(r - q) T} > d$. Pokud by tato podmínka nebyla splněna, vznikla by arbitrážní příležitost: například pokud $u \le e^{(r-q)T}$, prodali bychom akcii a investovali výnos do bezrizikového aktiva, čímž bychom dosáhli jistého zisku; obdobně pokud $d \ge e^{(r-q)T}$, půjčili bychom si a nakoupili akcii. % TODO: Zdroj
Jelikož se výdělky z obou portfólií rovnají, a dokážeme jednoznačně určit výdělek portfólia syntetického, musíme být schopni určit i hodnotu opčního portfolia.

\subsection{Risk-neutrální oceňování}

Rovnici \ref{eq_value} můžeme upravit na následující tvar:

\begin{align}
  \Delta S + \beta &= e^{-rT} \left( C_u \frac{e^{(r - q) T} -d}{u - d} + C_d \frac{u - e^{(r - q) T}}{u - d} \right) \\
                   &= e^{-rT} \left[ C_u \frac{e^{(r - q) T} -d}{u - d} + C_d \left( 1 - \frac{e^{(r - q) T} -d}{u - d} \right) \right].
\end{align}

Zlomek $\frac{e^{(r - q) T} - d}{u - d}$ vystupuje v roli váhy připadající $C_u$. Označíme-li

\begin{equation}
  \label{probability}
  p = \frac{e^{(r - q) T} - d}{u - d},
\end{equation}

potom můžeme říci, že $c$ vystupuje v rovnici jako pravděpodobnost růstu hodnoty ceny akcie, a můžeme hodnotu syntetické portfolia rozepsat jako:

\begin{equation}
  \Delta S + \beta = e^{-rT}\bigl(p \, C_u + (1-p) C_d\bigr). \label{eq_riskneutral}
\end{equation}

Aby $p$ skutečně představovalo pravděpodobnost, musí platit $0 < p < 1$. To je zaručeno arbitrážní podmínkou $u > e^{(r - q)T} > d$ — pokud by nebyla splněna, existoval by bezrizikový zisk. Veličina $p$ se nazývá \textbf{rizikově neutrální pravděpodobnost} %TODO: Citace — risk-neutral probability (Harrison & Kreps 1979, nebo textbook)
a odpovídá světu, v němž jsou všichni investoři lhostejní k riziku a očekávaný výnos všech aktiv je bezriziková sazba $r$. Nejedná se o skutečnou pravděpodobnost pohybu ceny akcie — tu neznáme a pro ocenění opce ji ani nepotřebujeme. Při pravděpodobnosti růstu ceny akcie $p$ je očekávaná cena akcie $S e^{(r-q)T}$. % TODO: Zdroj

Rovnice \eqref{eq_riskneutral} ukazuje, že cena opce je diskontovaná očekávaná výplata v rizikově neutrálním světě. Tento princip se nazývá \textbf{rizikově neutrální oceňování} (risk-neutral pricing) %TODO: Citace — risk-neutral pricing obecně
a je základem všech binomických modelů. 

\subsection{Binomické stromy s více kroky}

Zatím jsme počítali jen se stromy s pouze jedním krokem, čili stromy, kde množina vrcholů je $\{S, uS, dS\}$. Tento model bývá poněkud nepřesný, a abychom odvodili model vhodný pro americké opce, zavedeme model s více binomickými periodami.

Pro opce máme $T$ jako čas do data expirace, vyjádřený jako jedna perioda, nejčastěji rok. Pokud mluvíme o binomické periodě, myslíme délku času mezi dvěma pozorováními ceny akcie, které jsou v binomickém stromu spojeny hranou. Máme-li $n$ binomických period, délka binomické periody je $T/n$. 

Pokud chceme vytvořit binomický strom o $n$ binomických periodách, začneme vytvořím stromu začínajícím v počáteční ceně akcie $S$. Následně vytvoříme a s jejich rodičem, cenou akcie $S$ spojíme vrcholy $uS$ a $dS$. Z nově vzniklých vrcholů můžeme opět modelovat vzrůst či pád ceny. Jsme-li v libovolném listu $S_i$, můžeme z tohoto vrcholu pokračovat do vrcholů $u S_i$ a $d S_i$. Takto pokračujeme, dokud pro každý list ve stromě neplatí, že jeho vzdálenost od kořene (ceny akcie $S$) je $n$.

Nutno poukázat na následující fakt: Mějme binomický strom o vrcholech $\{S, uS, dS\}$. Ceny akcie vedoucí z listu $uS$ jsou ceny $u^2S$, $udS$ a z listu $dS$ vedou ceny $duS$ a $d^2S$. Jak vidíme, cena akci je stejná, pokud se propadne níž z vrcholu $uS$ a pokud stoupne z vrcholu $dS$. Tudíž, každé dva vrcholy, které mají stejného rodiče, mají jednoho společného potomka. Strom, pro který platí tento fakt, můžeme nazvat rekombinační.

\subsection{Zjištění ceny opce z binomického stromu}

Tímto jsme si dokázali odvodit vývoj ceny akcie pomocí binomického stromu. Nicméně, jak z tohoto odvodit cenu opce na danou akcii?

Hodnotu $p$ určíme z rovnice \ref{probability}. 

Začneme tím, že určíme cenu opce v listech binomického stromu. %TODO: Citace — vnitřní hodnota opce při expiraci
V listech opce vyprší — cena opce je tedy její vnitřní hodnota: $\max(S - K, 0)$ pro kupní opci a $\max(K - S, 0)$ pro prodejní opci. Následně ceny opcí v ostatních vrcholech určíme zpětnou indukcí: %TODO: Citace — zpětná indukce (backward induction)

\begin{equation}
  C = e^{-rT}\bigl(p C_u + (1-p) C_d\bigr).
\end{equation}

Takto postupujeme od listů až po kořen, čímž získáme současnou cenu opce.

Tímto jsme obecně zavedli ohodnocení binomickými stromy. 

\subsection{Dopočet hodnoty opce jako rovnice}

Mějme list $S_i$, pro který platí, že se cena akcie zvýšila právě $j$-krát. Potom každá cesta z kořene stromu do listu $S_i$ vzniká jako libovolná kombinace $j$ zvýšení cen a $n-j$ snížení cen, z čehož můžeme ukázat, že existuje právě $\binom{n}{j}$ cest vedoucích do vrcholu $S_i$.

Víme, že pro vrchol $S_i$ vzniká cena opce jako vážený součet potomků tohoto vrcholu, přičemž ceny potomků (nejsou-li listy stromu) také vznikají jako vážené součty jejich potomků. Z tohoto můžeme vyvodit, že cenu opce bez předčasného využití z binomického stromu můžeme vypočíst jako:

\begin{equation}
  C = e^{-rT} \sum_{i=1}^{n} p^i (1 - p)^{n - i} f(S u^i d^{n-i}),
\end{equation}

kde $f$ určuje výdělek opce při expiraci. Pro evropské kupní 


\subsection{Americké opce}

Zatímco u metod, které budou následovat, je oceňování amerických opcí složité a většinou vyžaduje zavedení dalšího aparátu, pro binomické stromy je toto velmi snadné: Pro každý vrchol, který není listem, spočítáme jak pokračovací hodnotu opce z navazujících vrcholů, tak i vnitřní hodnotu při okamžitém uplatnění. Cenou opce je větší z těchto dvou hodnot.

Pro kupní opci:

\begin{equation}
  C = \max\!\bigl(e^{-rT}(p C_u + (1-p) C_d),\, S - K\bigr),
\end{equation}

a pro prodejní opci:

\begin{equation}
  P = \max\!\bigl(e^{-rT}(p P_u + (1-p) P_d),\, K - S\bigr),
\end{equation}

kde $P_u$ je následující vrchol se zvýšenou cenou akcie a $P_d$ je následující vrchol se sníženou cenou akcie.

\subsection{Coxovy-Rossovy-Rubinsteinovy binomické stromy}

Binomické stromy se obecně liší ve třech způsobech: %TODO: Citace — varianty binomických stromů (CRR, Jarrow-Rudd, Tian, Leisen-Reimer apod.)

\begin{itemize}
  \item výše $u$ a $d$,
  \item pravděpodobnost růstu ceny akcie $p$ a pravděpodobnost poklesu ceny akcie $1 - p$
  \item centrování stromu.
\end{itemize}

Nejčastěji používaným modelem binomického stromu je Coxův-Rossův-Rubinsteinův binomický strom \citep[Kapitola~13]{McDonald2013} \citep[Kapitola~13]{Hull2021} a proto jej využijeme i v této práci.

Pro tento specifický model nastavíme:

\begin{gather}
  u = e^{(r - q)T + \sigma \sqrt{T}}, \\
  d = e^{(r - q)T - \sigma \sqrt{T}}.
\end{gather}

$e^{(r - q)T}$ je složka upravující cenu o vliv času skrze růst o bezrizikovou úrokovou míru a pokles o vyplácenou dividendu a $\sigma \sqrt{T}$ přidává nejistotu vývoje ceny. \citep[Kapitola~13]{McDonald2013}

\end{document}
